\documentclass[10pt,a4paper]{amsart}
\usepackage[margin=19mm]{geometry}
\usepackage{amsmath,amssymb,mathtools}
\usepackage[scr=rsfs,cal=euler]{mathalfa}
\usepackage{xfrac}
\usepackage[unicode,hidelinks,bookmarks=false]{hyperref}
\newcommand{\ie}{i.e.\ }
\newcommand{\cf}{cf.\ }
\newcommand{\resp}{respectively\ }
\newcommand{\Subsection}{\subsection}
\providecommand{\nobreakdash}{\nobreak\hbox{-}}
\setlength{\emergencystretch}{2em}
\multlinegap0pt
\usepackage{bm}
\usepackage{bbm}
\usepackage{dsfont}
\usepackage{xspace}
\usepackage{cancel}
\usepackage{mathtools}
\usepackage{amsthm}
\makeatletter
\@ifundefined{theorem}{\newtheorem{theorem}{Theorem}}{}
\@ifundefined{lemma}{\newtheorem{lemma}[theorem]{Lemma}}{}
\makeatother
\newcommand{\BD}{\mathcal {BD}}
\newcommand{\Q}{ \mathbb Q }
\newcommand{\w}{\tilde}
\title[Multiple q-zeta values and traces]{Multiple q-zeta values and traces}
\author{Zhenbo Qin}
\date{Source: arXiv:2505.14614v1 (CC BY 4.0)}
\begin{document}
\maketitle

\noindent\emph{Source and licence.} Multiple q-zeta values and traces. Version arXiv:2505.14614v1 (CC BY 4.0), distributed under the Creative Commons Attribution 4.0 International licence, as recorded in the arXiv licence metadata for this exact version and shown on the abstract page https://arxiv.org/abs/2505.14614v1. Full rights notice: licenses/arxiv-2505.14614-CC-BY-4.0.txt.\par\smallskip

\noindent\textit{Editorial scope.} Counted unit: the Lemma whose source label is \texttt{lemma\_t2t1yq} in the article (source0.tex lines 1044--1061, source label \texttt{lemma\_t2t1yq}), delivered together with the article's own complete original proof of it (source0.tex lines 1062--1313, one \texttt{proof} environment of 252 lines). No mathematics is written, repaired or re-derived: statement and proof are the article's own text line for line. Required context retained uncounted: sumofpowers (lines 853--866). Every definition, display and label of those lines is retained unchanged. Omitted from this package and not redistributed: 1--852, 867--1043, 1314--2317. Nothing inside a retained excerpt is omitted or altered: every retained body is delivered exactly as the article's lines are written, so no omission receipt of any kind accompanies this package. Citations: the retained excerpts carry no citation command at all, so no bibliography entry of the article is transcribed and no reference is redistributed. Reference adaptation: none was needed and none was made; every reference inside the delivered unit resolves inside it, so no unresolved reference or citation is produced. Printed slips kept literal: no printed word, symbol, hypothesis or number of the counted statement or of its proof was altered. Layout and class: the article's own class is \texttt{amsart} with packages including amsmath, amstext, amsbsy, amssymb, ytableau; the wrapper is the lane's pinned \texttt{amsart} editorial wrapper at 10pt on A4, which restates the theorem-like environments the retained text needs and adds the article's own macro definitions that the retained text needs (\texttt{\textbackslash BD}, \texttt{\textbackslash Q}, \texttt{\textbackslash w}), with extra wrapper packages (bm, bbm, dsfont, xspace, cancel, mathtools). The delivered numbering is the unit's own. Assets: no retained span contains a figure environment, a graphics-inclusion command, a PDF-page inclusion, a table or a TikZ picture; no image file is copied into this package, the package declares no asset dependency and no image file is redistributed with it. Licence: version arXiv:2505.14614v1 of this article is distributed under the Creative Commons Attribution 4.0 International licence, as recorded in the arXiv licence metadata for this exact version and shown on the abstract page https://arxiv.org/abs/2505.14614v1. A full rights notice is delivered at licenses/arxiv-2505.14614-CC-BY-4.0.txt. Characters: every retained excerpt is pure ASCII, so the delivered engine, XeLaTeX with its default Latin Modern text font, reports no missing character.
\begin{lemma}  \label{sumofpowers} 
Fix integers $n \ge 0$, $r \ge 1$ and $t_1, \ldots, t_{r} \ge 1$. Then, 
\begin{enumerate}
\item[{\rm (i)}]
$\displaystyle{\sum_{d_1, \ldots, d_r \ge 1, \, d_1+ \ldots + d_r \le n}
  \prod_{i=1}^{r} d_i^{t_i}}$
is a polynomial in $\Q[n] \cdot n$ of degree $\sum_{i=1}^{r} t_i+r$.

\item[{\rm (ii)}]
$\displaystyle{\sum_{d_1, \ldots, d_r \ge 1, \, d_1+ \ldots + d_r = n}
  \prod_{i=1}^{r} d_i^{t_i}}$
is a polynomial in $\Q[n] \cdot n$ of degree $\sum_{i=1}^{r} t_i+r-1$.
\end{enumerate}
\end{lemma}

\begin{lemma}  \label{lemma_t2t1yq} 
Fix integers $r, s \ge 1$ and $u_1, \ldots, u_{r+s}, 
t_1, \ldots, t_{r+s} \ge 0$. Then,
\begin{eqnarray}  \label{lemma_t2t1yq.0}
\sum_{\substack{n_1 > \ldots > n_r \ge 0, d_1, \ldots, d_r \ge 1\\
    n_{r+1} > \ldots > n_{r+s} \ge 1, d_{r+1}, \ldots, d_{r+s} \ge 1\\
    d_1+ \ldots + d_r = d_{r+1} + \ldots + d_{r+s}}}
  \prod_{i=1}^{r+s} n_i^{u_i}d_i^{t_i} \cdot q^{\sum_{i=1}^{r+s} n_i d_i}
\quad \in \quad \BD.
\end{eqnarray}
More precisely, the above summation is a finite $\Q$-linear combination of 
bi-brackets 
$$
\begin{bmatrix} t_1', \ldots, t_\ell'\\u_1', \ldots, u_\ell' \end{bmatrix}
$$
of weights at most $\sum_{i=1}^{r+s} (u_i + t_i+1)$. In addition, 
if $t_1, \ldots, t_{r+s} \ge 1$, then $t_1' \ge 2$. 
\end{lemma}

\begin{proof}
Use induction on $s$. 
Denote the summation in \eqref{lemma_t2t1yq.0} by $W_s$. 
For simplicity, we have suppressed the dependence of $W_s$ on 
$r, u_1, \ldots, u_{r+s}, t_1, \ldots, t_{r+s}$. 
The idea of the proof is to eliminate the parameters 
$d_{r+1}, \ldots, d_{r+s}$, via induction on $s$, from the condition 
$d_1+ \ldots + d_r = d_{r+1} + \ldots + d_{r+s}$ until it becomes void:
\begin{eqnarray}  \label{lemma_t2t1yq.100}
\w d_1+ \ldots + \w d_{\w r} = \w d_{\w r+1} + \ldots + \w d_{\w r+0}
\end{eqnarray}
($\w r$ may be different from $r$), 
which corresponds to a summation of the form $W_0$. 
During the elimination process, we may produce elements in $\BD$.

First of all, we rewrite $W_s$. Since $d_1+ \ldots + d_r 
- d_{r+1} - \ldots - d_{r+s-1} = d_{r+s} \ge 1$, 
$$
W_s = \sum_{\substack{n_1 > \ldots > n_r \ge 0, d_1, \ldots, d_r \ge 1\\
    n_{r+1} > \ldots > n_{r+s} \ge 1, d_{r+1}, \ldots, d_{r+s-1} \ge 1\\
    d_1+ \ldots + d_r > d_{r+1} + \ldots + d_{r+s-1}}}
  \prod_{i=1}^{r+s-1} n_i^{u_i}d_i^{t_i} \cdot n_{r+s}^{u_{r+s}}
  \left (\sum_{i=1}^r d_i - \sum_{i=1}^{s-1} d_{r+i} \right )^{t_{r+s}} 
$$
$$
\cdot q^{\sum_{i=1}^{r} (n_i + n_{r+s})d_i}
\cdot q^{\sum_{i=r+1}^{r+s-1} (n_i - n_{r+s})d_i}.
$$
So $W_s$ is a finite $\Q$-linear combination of expressions of the form
$$
\sum_{\substack{n_1 > \ldots > n_r \ge 0, d_1, \ldots, d_r \ge 1\\
  n_{r+1} > \ldots > n_{r+s} \ge 1, d_{r+1}, \ldots, d_{r+s-1} \ge 1\\
  d_1+ \ldots + d_r > d_{r+1} + \ldots + d_{r+s-1}}}
  \prod_{i=1}^{r+s-1} n_i^{u_i}d_i^{\w t_i} \cdot n_{r+s}^{u_{r+s}} 
  \cdot q^{\sum_{i=1}^{r} (n_i + n_{r+s})d_i}
  \cdot q^{\sum_{i=r+1}^{r+s-1} (n_i - n_{r+s})d_i}
$$
where $\w t_i \ge 0$ and $\sum_{i=1}^{r+s-1} \w t_i 
= \sum_{i=1}^{r+s} t_i$. 
After changes of variables (e.g, set $n = n_{r+s}$), 
$W_s$ is a finite $\Q$-linear combination of expressions of the form 
$$
\sum_{\substack{n_1 > \ldots > n_r \ge n \ge 1, d_1, \ldots, d_r \ge 1\\
    n_{r+1} > \ldots > n_{r+s-1} \ge 1, d_{r+1}, \ldots, d_{r+s-1} \ge 1\\
    d_1+ \ldots + d_r > d_{r+1} + \ldots + d_{r+s-1}}}
   \prod_{i=1}^{r} (n_i-n)^{u_i}d_i^{\w t_i} \cdot 
   \prod_{i=r+1}^{r+s-1} (n_i+n)^{u_i}d_i^{\w t_i} 
$$
$$
\cdot n^{u_{r+s}} \cdot q^{\sum_{i=1}^{r+s-1} n_i d_i}.
$$
which itself is a finite $\Q$-linear combination of expressions of the form
$$
\sum_{\substack{n_1 > \ldots > n_r \ge 1, d_1, \ldots, d_r \ge 1\\
    n_{r+1} > \ldots > n_{r+s-1} \ge 1, d_{r+1}, \ldots, d_{r+s-1} \ge 1\\
    d_1+ \ldots + d_r > d_{r+1} + \ldots + d_{r+s-1}}}
   \prod_{i=1}^{r+s-1} n_i^{\w u_i}d_i^{\w t_i}  
   \cdot q^{\sum_{i=1}^{r+s-1} n_i d_i} \cdot \sum_{n=1}^{n_r} n^a
$$
where $\w u_i, a \ge 0$ and $\sum_{i=1}^{r+s-1} \w u_i +a 
= \sum_{i=1}^{r+s} u_i$. By Lemma~\ref{sumofpowers}~(i), 
$W_s$ is a finite $\Q$-linear combination of expressions of the form
\begin{eqnarray}  \label{lemma_t2t1yq.1}
\sum_{\substack{n_1 > \ldots > n_r \ge 1, d_1, \ldots, d_r \ge 1\\
    n_{r+1} > \ldots > n_{r+s-1} \ge 1, d_{r+1}, \ldots, d_{r+s-1} \ge 1\\
    d_1+ \ldots + d_r > d_{r+1} + \ldots + d_{r+s-1}}}
   \prod_{i=1}^{r+s-1} n_i^{\hat u_i}d_i^{\w t_i}  
   \cdot q^{\sum_{i=1}^{r+s-1} n_i d_i}
\end{eqnarray}
where $\hat u_i, \w t_i \ge 0$, $\hat u_r \ge 1$ and 
$\sum_{i=1}^{r+s-1} \hat u_i \le 
\sum_{i=1}^{r+s-1} \w u_i +a + 1 = \sum_{i=1}^{r+s} u_i + 1$. So
\begin{eqnarray}  \label{lemma_t2t1yq.2}
\sum_{i=1}^{r+s-1} (\hat u_i + \w t_i + 1)   
\le \sum_{i=1}^{r+s} (u_i + t_i + 1).
\end{eqnarray}
Moreover,
if $t_1, \ldots, t_{r+s} \ge 1$, then $\w t_1, \ldots, \w t_{r+s-1} \ge 1$. 
%Similarly, if $u_1, \ldots, u_{r+s} \ge 1$, 
%then $\hat u_1, \ldots, \hat u_{r+s-1} \ge 1$.

When $s=1$, we see from \eqref{lemma_t2t1yq.1} and \eqref{lemma_t2t1yq.2}
that $W_1$ is a finite $\Q$-linear combination of expressions of the form
$$
\sum_{n_1 > \ldots > n_r \ge 1, d_1, \ldots, d_r \ge 1}
   \prod_{i=1}^{r} n_i^{\hat u_i}d_i^{\w t_i}  
   \cdot q^{\sum_{i=1}^{r} n_i d_i}
= \prod_{i=1}^{r} (\hat u_i! \cdot \w t_i!) \cdot 
   \begin{bmatrix} \w t_1 + 1, \ldots, \w t_r + 1\\
     \hat u_1, \ldots, \hat u_r 
   \end{bmatrix}
$$
where $\sum_{i=1}^{r} (\hat u_i + \w t_i + 1)   
\le \sum_{i=1}^{r+1} (u_i + t_i + 1)$. Moreover,
if $t_1, \ldots, t_{r+1} \ge 1$, then $\w t_1, \ldots, \w t_{r} \ge 1$. 
%Similarly, if $u_1, \ldots, u_{r+1} \ge 1$, 
%then $\hat u_1, \ldots, \hat u_{r} \ge 1$. 
Hence the lemma holds for $s=1$.

Next, fix $s \ge 2$. Assume that the lemma holds for all the $W_\ell$'s 
(with various integers $r \ge 1, u_1, \ldots, u_{r+\ell}, 
t_1, \ldots, t_{r+\ell} \ge 0$) whenever $1 \le \ell \le s-1$. 
We will prove that the lemma holds for $W_s$ as well.
Rewrite \eqref{lemma_t2t1yq.1} as
\begin{eqnarray}  \label{lemma_t2t1yq.3}
\sum_{\substack{n_1 > \ldots > n_r \ge 1, d_1, \ldots, d_r \ge 1\\
    n_{r+1} > \ldots > n_{r+s-1} \ge 1, d_{r+1}, \ldots, d_{r+s-1} \ge 1}}
- \sum_{\substack{n_1 > \ldots > n_r \ge 1, d_1, \ldots, d_r \ge 1\\
    n_{r+1} > \ldots > n_{r+s-1} \ge 1, d_{r+1}, \ldots, d_{r+s-1} \ge 1\\
    d_1+ \ldots + d_r = d_{r+1} + \ldots + d_{r+s-1}}}
\end{eqnarray}
\begin{eqnarray}  \label{lemma_t2t1yq.4}
- \sum_{\substack{n_1 > \ldots > n_r \ge 1, d_1, \ldots, d_r \ge 1\\
    n_{r+1} > \ldots > n_{r+s-1} \ge 1, d_{r+1}, \ldots, d_{r+s-1} \ge 1\\
    d_1+ \ldots + d_r < d_{r+1} + \ldots + d_{r+s-1}}}.
\end{eqnarray}
The first summation in \eqref{lemma_t2t1yq.3} is equal to
$$
\prod_{i=1}^{r+s-1} (\hat u_i! \cdot \w t_i!) \cdot 
   \begin{bmatrix} \w t_1 + 1, \ldots, \w t_r + 1\\
     \hat u_1, \ldots, \hat u_r 
   \end{bmatrix}
\cdot \begin{bmatrix} \w t_{r+1} + 1, \ldots, \w t_{r+s-1} + 1\\
     \hat u_{r+1}, \ldots, \hat u_{r+s-1}
   \end{bmatrix}
\in \BD.
$$
Next, we deal with \eqref{lemma_t2t1yq.4}. 
Note that \eqref{lemma_t2t1yq.4} is equal to
$$
- \sum_{\substack{n_1 > \ldots > n_r \ge 1, d_1, \ldots, d_r, d \ge 1\\
    n_{r+1} > \ldots > n_{r+s-1} \ge 1, d_{r+1}, \ldots, d_{r+s-1} \ge 1\\
    d_1+ \ldots + d_r + d = d_{r+1} + \ldots + d_{r+s-1}}}
    \prod_{i=1}^{r+s-1} n_i^{\hat u_i}d_i^{\w t_i}  
    \cdot q^{\sum_{i=1}^{r+s-1} n_i d_i}.
$$
Since $d_1+ \ldots + d_r + d - d_{r+1} - \ldots - d_{r+s-2} 
= d_{r+s-1} \ge 1$, \eqref{lemma_t2t1yq.4} is equal to
$$
- \sum_{\substack{n_1 > \ldots > n_r \ge 1, d_1, \ldots, d_r, d \ge 1\\
    n_{r+1} > \ldots > n_{r+s-1} \ge 1, d_{r+1}, \ldots, d_{r+s-2} \ge 1\\
    d_1+ \ldots + d_r + d > d_{r+1} + \ldots + d_{r+s-2}}}
    \prod_{i=1}^{r+s-2} n_i^{\hat u_i}d_i^{\w t_i} \cdot 
    n_{r+s-1}^{\hat u_{r+s-1}} \left (\prod_{i=1}^{r} 
    d_i + d - \prod_{i=1}^{s-2} d_{r+i} \right )^{\w t_{r+s-1}}
$$
$$
\cdot q^{\sum_{i=1}^{r+s-2} n_i d_i}
\cdot q^{n_{r+s-1} (d_1+ \ldots + d_r + d - d_{r+1} - \ldots - d_{r+s-2})}
$$
which is a finite $\Q$-linear combination of expressions of the form
$$
\sum_{\substack{n_1 > \ldots > n_r > n \ge 1, d_1, \ldots, d_r, d \ge 1\\
    n_{r+1} > \ldots > n_{r+s-2} \ge 1, d_{r+1}, \ldots, d_{r+s-2} \ge 1\\
    d_1+ \ldots + d_r + d > d_{r+1} + \ldots + d_{r+s-2}}}
    \prod_{i=1}^{r} (n_i-n)^{\hat u_i}d_i^{\hat t_i} \cdot 
    \prod_{i=r+1}^{r+s-2} (n_i+n)^{\hat u_i}d_i^{\hat t_i} 
$$
$$
\cdot n^{\hat u_{r+s-1}} d^{\hat t_{r+s-1}}
\cdot q^{\sum_{i=1}^{r+s-2} n_i d_i} \cdot q^{nd}
$$
where $\hat t_i \ge 0$ ({\it it is possible that $\hat t_{r+s-1}=0$ 
even if $t_1, \ldots, t_{r+s} \ge 1$}) and 
$$
  \sum_{i=1}^{r+s-1} \hat t_i = \sum_{i=1}^{r+s-1} \w t_i 
= \sum_{i=1}^{r+s} t_i. 
$$
It follows that \eqref{lemma_t2t1yq.4} 
is a finite $\Q$-linear combination of expressions of the form
\begin{eqnarray}  \label{lemma_t2t1yq.6}
\sum_{\substack{n_1 > \ldots > n_r > n \ge 1, d_1, \ldots, d_r, d \ge 1\\
    n_{r+1} > \ldots > n_{r+s-2} \ge 1, d_{r+1}, \ldots, d_{r+s-2} \ge 1\\
    d_1+ \ldots + d_r + d > d_{r+1} + \ldots + d_{r+s-2}}}
    \prod_{i=1}^{r+s-2} n_i^{\overline u_i}d_i^{\hat t_i} \cdot 
    n^{\overline u_{r+s-1}} d^{\hat t_{r+s-1}}
    \cdot q^{\sum_{i=1}^{r+s-2} n_i d_i + nd}
\end{eqnarray} 
where $\overline u_i \ge 0$ and $\sum_{i=1}^{r+s-1} \overline u_i 
= \sum_{i=1}^{r+s-1} \hat u_i$. Note that 
$$
    \sum_{i=1}^{r+s-1} (\overline u_i + \hat t_i + 1)  
= \sum_{i=1}^{r+s-1} (\hat u_i + \w t_i + 1)    
\le \sum_{i=1}^{r+s} (u_i + t_i + 1)
$$
by \eqref{lemma_t2t1yq.2}, and that if $t_1, \ldots, t_{r+1} \ge 1$, 
then $\hat t_1, \ldots, \hat t_{r} \ge 1$. Moreover, 
\eqref{lemma_t2t1yq.6} is of the form \eqref{lemma_t2t1yq.1} 
with the index $s$ there being replaced by $s-1$,
and we repeat the above argument beginning at line \eqref{lemma_t2t1yq.3}.

Now we deal with the second summation in \eqref{lemma_t2t1yq.3} 
which is equal to
\begin{eqnarray}  \label{lemma_t2t1yq.5}
\sum_{\substack{n_1 > \ldots > n_r \ge 0, d_1, \ldots, d_r \ge 1\\
    n_{r+1} > \ldots > n_{r+s-1} \ge 1, d_{r+1}, \ldots, d_{r+s-1} \ge 1\\
    d_1+ \ldots + d_r = d_{r+1} + \ldots + d_{r+s-1}}} 
    \prod_{i=1}^{r+s-1} n_i^{\hat u_i}d_i^{\w t_i}  
    \cdot q^{\sum_{i=1}^{r+s-1} n_i d_i}
\end{eqnarray}
$$
- 0^a \cdot 
  \sum_{\substack{n_1 > \ldots > n_{r-1} \ge 1, d_1, \ldots, d_{r-1} \ge 1\\
    n_{r} > \ldots > n_{r+s-2} \ge 1, d_{r}, \ldots, d_{r+s-2} \ge 1\\
    d_1+ \ldots + d_{r-1} < d_{r} + \ldots + d_{r+s-2}}}
    \prod_{i=1}^{r+s-2} n_i^{\overline u_i} d_i^{\overline t_i}  
    \cdot q^{\sum_{i=1}^{r+s-2} n_i d_i} \cdot  
    \left (\sum_{i=0}^{s-2} d_{r+i} - \sum_{i=1}^{r-1} d_i \right )^b
$$
where $a = \hat u_r \ge 0$ (we set $0^0 =1$), $b = \w t_r$, 
we have renamed $\hat u_i$ and $\w t_i$ as 
$\overline u_i$ and $\overline t_i$ 
respectively when $1 \le i \le r-1$, and we have renamed 
$n_{r+i}, d_{r+i}, \hat u_{r+i}, \w t_{r+i}$ as $n_{r+i-1}, d_{r+i-1}, 
\overline u_{r+i-1}, \overline t_{r+i-1}$ respectively 
when $1 \le i \le s-1$. So 
$$
    \sum_{i=1}^{r+s-2} (\overline u_i + \overline t_i + 1) + b 
\le \sum_{i=1}^{r+s-1} (\hat u_i + \w t_i + 1) - 1
< \sum_{i=1}^{r+s} (u_i + t_i + 1)
$$
by \eqref{lemma_t2t1yq.2}. Since line \eqref{lemma_t2t1yq.5} is 
of the form $W_{s-1}$, we deduce from induction that 
line \eqref{lemma_t2t1yq.5} is 
a finite $\Q$-linear combination of bi-brackets of weights at most 
$$
  \sum_{i=1}^{r+s-1} (\hat u_i + \w t_i + 1)
\le \sum_{i=1}^{r+s} (u_i + t_i + 1).
$$
The line below line \eqref{lemma_t2t1yq.5} is a finite $\Q$-linear combination of the expressions 
$$
0^a \cdot 
\sum_{\substack{n_1 > \ldots > n_{r-1} \ge 1, d_1, \ldots, d_{r-1} \ge 1\\
    n_{r} > \ldots > n_{r+s-2} \ge 1, d_{r}, \ldots, d_{r+s-2} \ge 1\\
    d_1+ \ldots + d_{r-1} < d_{r} + \ldots + d_{r+s-2}}}
    \prod_{i=1}^{r+s-2} n_i^{\overline u_i} d_i^{\hat t_i}  
    \cdot q^{\sum_{i=1}^{r+s-2} n_i d_i}
$$
where $\hat t_i \ge 0$ and $\sum_{i=1}^{r+s-2} \hat t_i 
= \sum_{i=1}^{r+s-2} t_i + b$. The preceding line is either $0$ or 
is of the form \eqref{lemma_t2t1yq.4} with the index $r$ there 
being replaced by $r-1$, and so we repeat the above argument for 
\eqref{lemma_t2t1yq.4}.

In summary, we conclude that $W_s$ is a finite $\Q$-linear combination of 
bi-brackets 
$$
\begin{bmatrix} t_1', \ldots, t_\ell'\\u_1', \ldots, u_\ell' \end{bmatrix}
$$
of weights at most $\sum_{i=1}^{r+s} (u_i + t_i+1)$. In addition, 
if $t_1, \ldots, t_{r+s} \ge 1$, then $t_1' \ge 2$. 
\end{proof}

\end{document}
